\documentclass[11pt,a4paper]{amsart}

\usepackage[T1]{fontenc}
\usepackage{amsmath,amssymb,amsthm,mathtools}
\usepackage{enumitem}
\usepackage{microtype}
\usepackage[margin=2.7cm]{geometry}
\usepackage{xcolor}
\definecolor{linkcol}{rgb}{0.10,0.20,0.50}
\usepackage[colorlinks=true,linkcolor=linkcol,citecolor=linkcol,urlcolor=linkcol]{hyperref}

\theoremstyle{plain}
\newtheorem{theorem}{Theorem}[section]
\newtheorem{proposition}[theorem]{Proposition}
\newtheorem{lemma}[theorem]{Lemma}
\newtheorem{corollary}[theorem]{Corollary}
\theoremstyle{definition}
\newtheorem{definition}[theorem]{Definition}
\newtheorem{example}[theorem]{Example}
\newtheorem{remark}[theorem]{Remark}

\newcommand{\RR}{\mathbb R}
\newcommand{\QQ}{\mathbb Q}
\newcommand{\CC}{\mathbb C}
\newcommand{\NN}{\mathbb N}
\newcommand{\one}{\mathbf 1}
\newcommand{\Irr}{\operatorname{Irr}}
\newcommand{\Hom}{\operatorname{Hom}}
\newcommand{\im}{\operatorname{im}}
\newcommand{\indm}{\operatorname{ind}_{-}}
\newcommand{\corank}{\operatorname{corank}}
\newcommand{\rank}{\operatorname{rank}}
\newcommand{\Span}{\operatorname{span}}
\newcommand{\Sh}{\operatorname{Sh}}
\newcommand{\Vect}{\operatorname{Vect}}
\newcommand{\Hess}{\operatorname{Hess}}
\newcommand{\Def}{\operatorname{Def}}
\newcommand{\Tcal}{\mathcal T}
\newcommand{\Pcal}{\mathcal P}
\newcommand{\Mcal}{\mathcal M}
\newcommand{\Ecal}{\mathcal E}

\begin{document}

\title[The Stokes--microlocal Witt law]{The polarized Stokes--microlocal Witt law}
\author{M.~Haufschild}
\address{Independent researcher}
\email{haufschildmichael@gmail.com}
\date{4 August 2026}

\subjclass[2020]{Primary 34M40, 32S60; Secondary 14D24, 15A18, 18F20, 16G20}
\keywords{Stokes structure, irregular Riemann--Hilbert correspondence, microsupport,
vanishing cycle, characteristic cycle, Rees construction, wild character variety, Witt index,
Puiseux contact, quasi-Hamiltonian reduction}

\begin{abstract}
For a finite Galois-stable Puiseux spectrum, highest coefficient difference determines mixed Hom
irregularity.  We prove that its threshold partitions form a canonical dendrogram and obtain an
exact coarea decomposition of the full polarized Stokes index.  Every threshold term is a positive
square; physical rank is the sole potentially negative channel.

We construct the finite cellular Stokes category on the boundary circle.  Microlocalization is a
quotient functor, nearby-to-vanishing is an exact quotient, and the cellular characteristic cycle
of a degree-$k$ exponential Hom germ has $2k$ oriented rays of unit multiplicity, exactly $k$ in
the ordered decay orientation.  Galois-normalized directed microstalk mass equals mixed Hom
irregularity.  Thus the coarea decomposition is also a categorical microlocal index formula.

After physical-rank fusion over an arbitrary marked curve, all local channels remain positive and
topology changes only one common-rank square.  The global form has Witt index at most one for every
genus, puncture count, ramification pattern, and nonnegative multi-Rees specialization.  Capacity
completion gives the sharp finite/affine/Lorentz trichotomy.  Corona reduction is an exact positive
Schur complement; its residual is the quotient in an exact boundary--total--relative-vanishing
sequence, and successive eliminations are order-independent.  An explicit three-term deformation
complex derives the smooth-reduction dimension and Euler formulas by rank--nullity.  These results
give a mixed, polarized refinement of scalar irregular rigidity formulae.
\end{abstract}

\maketitle

\section{Introduction}

Irregular geometry combines two kinds of gluing.  Locally, a ramified exponential spectrum is
split into branches on a finite cover and reglued by formal monodromy.  Globally, local Stokes
spaces are fused after their physical ranks have been identified.  Both operations obscure a
simple fact proved here: all local deformation channels form a positive coarea energy, while every
possible negative contribution factors through one scalar, total physical rank.

The scalar ancestry is important.  Hiroe proves, under multiplicity-free, regularity, and
spectral-curve hypotheses, that an irregular rigidity index is the Euler characteristic of a
normalized spectral curve \cite{Hiroe2019}; related local formulas involve Milnor numbers and
intersection multiplicities \cite{Hiroe2017}.  We neither rename nor reprove that theorem.  Our
invariant is the whole polarization on independent orbit charges.  Its coefficients are all mixed
$\Irr\Hom$ terms, and its conclusion is an inertia theorem rather than one scalar evaluation.

Let $w^N=z$ be a common local cover and let
\[
                     q_a(w)=\sum_{k=1}^{K}c_{a,k}w^{-k}
\]
run through a finite $\mu_N$-stable set $A$ of distinct exponential factors.  For $a\ne b$, let
$k(a,b)$ be the largest exponent whose coefficients differ and put $h(a,b)=k(a,b)/N$; set
$h(a,a)=0$.  The elementary inequality
\[
                     h(a,b)\leq\max\{h(a,c),h(c,b)\}
\]
forces the threshold partitions of $A$ to be nested.  If $H$ lies above all contacts, the resulting
finite dendrogram has blocks $A_e$ and positive persistence lengths $\ell_e$.  If
$I=A/\mu_N$, if $x_i$ is an orbit charge, and if $y_a=x_{[a]}$, our local theorem is
\begin{equation}\label{eq:introduction-law}
 \boxed{
 Q_A(x)=\sum_{i\in I}x_i^2+
        \sum_e\ell_e\left(\sum_{a\in A_e}y_a\right)^2
        -(H-1)\left(\sum_{a\in A}y_a\right)^2.}              \tag{W}
\end{equation}
No tree, cut family, or contact matrix occurs among the hypotheses.  Each is reconstructed from
coefficients.  Branch blocks need not be unions of Galois orbits; the vector attached to $A_e$ is
therefore the exact incidence-count vector
\[
                       u_{e,i}=|A_e\cap i|,
\]
not an assumed orbit indicator.  This distinction is essential for ramified data.

Formula \eqref{eq:introduction-law} has four further realizations.
\begin{enumerate}[label=(\roman*),leftmargin=2.2em]
\item The linear form $\nu_e(x)=\sum_i u_{e,i}x_i$ is the principal conormal of the cluster
hyperplane $\nu_e=0$.
\item The associated-graded coordinate is $\sqrt{\ell_e}\,\nu_e$.  Independent nonnegative Rees
parameters scale these coordinates without changing the rank channel.
\item A degree-$k$ exponential Hom germ has directed microstalk mass $k$.  After orbit summation
and normalization by $N$, these masses are exactly the coefficients
$s_{ij}=\Irr\Hom(q_i,q_j)$ in $Q_A$.
\item Eliminating a positive corona block is exact Schur complementation.  The discarded square is
the positive norm of a relative-vanishing quotient coordinate.
\end{enumerate}

For a genus-$g$ curve with a nonempty set of $m$ marked points, identify all local physical ranks
and write their common value as $M$.  The fused form is
\begin{equation}\label{eq:introduction-fusion}
                   Q_{X,D}=\sum_{a=1}^{m}Q_{A_a}+2(1-g-m)M^2.
\end{equation}
Substituting \eqref{eq:introduction-law} leaves a positive-definite sum plus one multiple of
$M^2$.  Hence
\begin{equation}\label{eq:introduction-index}
             Q_{X,D}>0\ \hbox{on }\ker M\setminus\{0\},\qquad
             \indm(Q_{X,D})+\dim\ker(Q_{X,D})\leq1.          \tag{G}
\end{equation}
This is the global Stokes--microlocal Witt law.  Capacity completion replaces the displayed
rank coefficient by the intrinsic constrained minimum at $M=1$ and yields the exact
finite/affine/Lorentz trichotomy.  Every handle changes only this coefficient, by $-2$.

The paper works entirely with finite Puiseux, cellular Stokes, and tangent-reduction objects.  The
analytic irregular Riemann--Hilbert correspondence supplies the surrounding equivalence of
languages \cite{SabbahStokes,DKRH}; it is not a hidden proof input.  Section~\ref{sec:transport}
proves the precise transport statement needed here: any exact realization equivalence carrying
the two arc-to-wall maps to the cellular germ carries the quotient microstalk and all numerical
terms in (W).  Thus no analytic comparison is inferred from equality of numbers.

Sections~\ref{sec:coefficients}--\ref{sec:global} prove the coefficient and global laws.
Sections~\ref{sec:cellular}--\ref{sec:transport} construct the categorical microlocal side.
Sections~\ref{sec:rees}--\ref{sec:corona} treat Rees, capacity, and corona reduction.
Section~\ref{sec:deformation} derives the deformation dimension.  The final sections give marking
legs, spectral consequences, novelty boundaries, and falsifiers.

\section{Finite Puiseux spectra and their canonical cuts}\label{sec:coefficients}

\begin{definition}[Equivariant finite Puiseux spectrum]\label{def:spectrum}
Fix $N,K\geq1$ and a primitive $N$th root $\zeta$.  An equivariant finite Puiseux spectrum is a
finite set $A$, a permutation $\sigma:A\to A$ with $\sigma^N=1$, and coefficients
$c_{a,k}\in\CC$ such that
\[
                  c_{\sigma a,k}=\zeta^{-k}c_{a,k}
                  \qquad(a\in A,\ 1\leq k\leq K).
\]
We require the coefficient vectors $(c_{a,1},\ldots,c_{a,K})$ to be pairwise distinct.  The
associated exponential factor is $q_a(w)=\sum_kc_{a,k}w^{-k}$.
\end{definition}

For $a\ne b$, define
\[
 k(a,b)=\max\{k:c_{a,k}\ne c_{b,k}\},\qquad h(a,b)=k(a,b)/N,
\]
and put $k(a,a)=h(a,a)=0$.

\begin{lemma}[Highest-difference three-point law]\label{lem:three-point}
For all $a,b,c\in A$,
\begin{equation}\label{eq:three-point}
                         k(a,b)\leq\max\{k(a,c),k(c,b)\}.
\end{equation}
Consequently the largest of $k(a,b),k(a,c),k(b,c)$ occurs at least twice.
\end{lemma}

\begin{proof}
Above the maximum on the right, coefficients of $a$ and $c$ agree and coefficients of $c$ and $b$
agree.  Hence coefficients of $a$ and $b$ agree.  This proves \eqref{eq:three-point}.  Applying the
same inequality after permuting $a,b,c$ excludes a unique largest value.
\end{proof}

\section{Capacity completion and the intrinsic phase}\label{sec:capacity}

Write the positive local form as $P_A(x)=x^TPx$ and physical rank as $M(x)=r^Tx$, where
$r_i=|i|$.  Thus
\[
                         Q_A(x)=x^TPx-\lambda(r^Tx)^2,
                         \qquad\lambda=H-1,
\]
and $P$ is positive definite.  Define
\begin{equation}\label{eq:capacity-definitions}
                  c=r^TP^{-1}r>0,qquad \Omega=c^{-1}-\lambda.
\end{equation}

\begin{theorem}[Capacity normal form]\label{thm:capacity-normal}
For $m=r^Tx$, put
\begin{equation}\label{eq:centered-vector}
                       z=x-\frac{m}{c}P^{-1}r.
\end{equation}
Then $r^Tz=0$, this decomposition is unique, and
\begin{equation}\label{eq:capacity-normal}
                       \boxed{Q_A(x)=z^TPz+\Omega m^2.}       \tag{K}
\end{equation}
Consequently:
\begin{enumerate}[label=(\roman*),leftmargin=2.2em]
\item if $\Omega>0$, then $Q_A$ is positive definite;
\item if $\Omega=0$, then $Q_A$ is positive semidefinite with kernel
$\RR P^{-1}r$;
\item if $\Omega<0$, then $Q_A$ is nonsingular with exactly one negative eigenvalue.
\end{enumerate}
Moreover
\begin{equation}\label{eq:determinant-wall}
                         \det Q_A=\det(P)c\Omega.
\end{equation}
\end{theorem}

\begin{proof}
Since $r^TP^{-1}r=c$, equation \eqref{eq:centered-vector} gives $r^Tz=0$.  Applying $r^T$ proves
uniqueness.  Orthogonality $z^TP(P^{-1}r)=z^Tr=0$ gives
\[
 x^TPx-\lambda m^2=z^TPz+(c^{-1}-\lambda)m^2.
\]
This proves \eqref{eq:capacity-normal}.  Its orthogonal positive-space-plus-one-line decomposition
gives the trichotomy.  The matrix determinant lemma gives
\[
 \det(P-\lambda rr^T)=\det(P)(1-\lambda c)=\det(P)c(c^{-1}-\lambda).
\]
\end{proof}

\begin{proposition}[Root-height gauge invariance]\label{prop:root-gauge}
Increasing $H$ by $t>0$ adds the root block $A$ with persistence $t$.  Hence
\[
 P_t=P+trr^T,\qquad \lambda_t=\lambda+t,\qquad
 c_t=\frac{c}{1+tc}.
\]
In particular $c_t^{-1}-\lambda_t=\Omega$ and
$\det(P_t)c_t=\det(P)c$.
\end{proposition}

\begin{proof}
The charge of the full block is $M=r^Tx$, giving the first two formulas.  Sherman--Morrison gives
the formula for $c_t$; the two invariants follow immediately, using the determinant lemma for the
second.
\end{proof}

For every puncture let $c_a,\Omega_a$ be its local capacity data and put
$\delta_a=2-\Omega_a$.

\begin{theorem}[Euler--capacity refinement]\label{thm:euler-capacity}
Every $x=(x_a)\in V_{X,D}$ has a unique decomposition
$x_a=z_a+(M/c_a)P_a^{-1}r_a$, with $r_a^Tz_a=0$, and
\begin{equation}\label{eq:euler-capacity}
 \boxed{
 Q_{X,D}(x)=\sum_a z_a^TP_az_a+\Omega_{X,D}M^2,
 \qquad
 \Omega_{X,D}=\chi(X)+\sum_a(\Omega_a-2)
              =\chi(X)-\sum_a\delta_a.}                    \tag{EC}
\end{equation}
The scalar $\Omega_{X,D}$ is the constrained minimum of $Q_{X,D}$ at $M=1$.  Its sign gives the
positive-definite, corank-one affine, and Lorentz phases of the global form.
\end{theorem}

\begin{proof}
Apply Theorem~\ref{thm:capacity-normal} at every puncture and insert the result into
\eqref{eq:fused-form}.  The coefficient of $M^2$ is
\[
 \sum_a\Omega_a+2(1-g-m)=2-2g+\sum_a(\Omega_a-2).
\]
This is \eqref{eq:euler-capacity}.  The positive sum is minimized uniquely at $z_a=0$ for all
$a$, proving the variational claim and trichotomy.
\end{proof}

\section{Corona elimination as relative vanishing}\label{sec:corona}

Let $X,Y$ be finite-dimensional real vector spaces with coordinates, and let
\[
             \mathcal Q(x,y)=
             \begin{pmatrix}x\\y\end{pmatrix}^{T}
             \begin{pmatrix}A&B\\B^T&D\end{pmatrix}
             \begin{pmatrix}x\\y\end{pmatrix},
             \qquad D>0.
\]
Define the harmonic extension and relative-vanishing maps
\begin{equation}\label{eq:corona-maps}
 h(x)=(x,-D^{-1}B^Tx),\qquad
 v(x,y)=y+D^{-1}B^Tx.
\end{equation}

\begin{theorem}[Exact corona--vanishing law]\label{thm:corona}
There is a short exact sequence
\begin{equation}\label{eq:corona-exact}
              0\longrightarrow X\xrightarrow{h}X\oplus Y
              \xrightarrow{v}Y\longrightarrow0.             \tag{V}
\end{equation}
Moreover
\begin{equation}\label{eq:schur-completion}
 \boxed{
 \mathcal Q(x,y)=v(x,y)^TDv(x,y)+
 x^T(A-BD^{-1}B^T)x.}                                      \tag{S}
\end{equation}
Thus the Schur complement is the unique exact effective boundary energy; equality with it occurs
exactly on the harmonic graph $\im h$.  The quotient
$(X\oplus Y)/\im h\cong Y$ is represented by $v$, and the energy lost under elimination is its
positive $D$-norm.
\end{theorem}

\begin{proof}
The map $h$ is injective and $v$ is surjective since $v(0,y)=y$.  Also $v\circ h=0$.  Conversely,
$v(x,y)=0$ implies $y=-D^{-1}B^Tx$, so $(x,y)=h(x)$.  This proves exactness.
Expanding the first term on the right of \eqref{eq:schur-completion} cancels the two Schur terms and
recovers $x^TAx+2x^TBy+y^TDy$.  Positivity of $D$ proves the minimization and equality statements.
\end{proof}

\begin{corollary}[Order independence]\label{cor:corona-order}
Successive eliminations of positive hidden corona blocks are independent of elimination order.
\end{corollary}

\begin{proof}
Each elimination is exact minimization over one hidden block.  Iterated minima over a finite
product do not depend on the order.  Equivalently, an exact effective energy is unique: if two
functions are attained minima and lower bounds for the same full energy, each bounds the other.
\end{proof}

Theorem~\ref{thm:corona} resolves the terminology precisely.  A corona residual is not merely
called a vanishing channel because its square is positive: it is the coordinate of the explicit
quotient in \eqref{eq:corona-exact}.  In a cellular realization, Theorem~\ref{thm:transport}
carries this quotient to the corresponding relative vanishing object.

\section{The Stokes deformation complex}\label{sec:deformation}

We derive the dimension identity from tangent data.  Fix $p$ punctures and nonnegative integers
\[
 g,\quad G=\dim\mathfrak g,\quad z=\dim Z(\mathfrak g),\quad
 h_a,\ s_a,\ k_a\quad(1\leq a\leq p),
\]
where $h_a$ is formal-centralizer dimension, $s_a$ is directed Stokes-root dimension, and $k_a$
is prescribed conjugacy-class dimension.  Set
\begin{align}\label{eq:ambient-reduction-counts}
 A_1&=2gG+\sum_a(G+h_a+s_a+k_a),\\
 R&=G-z+\sum_ah_a.\notag
\end{align}

At a smooth stable reduction point, let
\[
 d_{\mathrm{act}}:\CC^R\longrightarrow\CC^{A_1},\qquad
 d_\mu:\CC^{A_1}\longrightarrow\CC^R
\]
be the effective gauge derivative and moment derivative.  Smooth stability means
$d_{\mathrm{act}}$ is injective, $d_\mu$ is surjective, and moment invariance gives
$d_\mu d_{\mathrm{act}}=0$.

\begin{theorem}[Source-derived tangent dimension]\label{thm:tangent-dimension}
The reduced tangent space
\[
                         T=\ker d_\mu/\im d_{\mathrm{act}}
\]
has dimension
\begin{equation}\label{eq:tangent-dimension}
 \boxed{
 \dim T=2z+(2g-2+p)G+\sum_a(s_a+k_a-h_a).}                  \tag{D}
\end{equation}
\end{theorem}

\begin{proof}
Surjectivity gives $\dim\ker d_\mu=A_1-R$; injectivity gives
$\dim\im d_{\mathrm{act}}=R$.  Hence $\dim T=A_1-2R$.  Substitute
\eqref{eq:ambient-reduction-counts} and collect terms.
\end{proof}

Add a central summand in degrees zero and two and form
\begin{equation}\label{eq:deformation-complex}
 \CC^z\oplus\CC^R\xrightarrow{d_0}\CC^{A_1}
 \xrightarrow{d_1}\CC^R\oplus\CC^z,
 \qquad
 d_0(c,u)=d_{\mathrm{act}}u,\quad d_1(v)=(d_\mu v,0).
\end{equation}

\begin{theorem}[Self-dual Euler count]\label{thm:deformation-euler}
Sequence \eqref{eq:deformation-complex} is a complex and
\[
 \dim H^0=z,\qquad H^1\cong T,\qquad \dim H^2=z.
\]
Therefore
\begin{equation}\label{eq:deformation-euler}
 \boxed{
 \chi(\Def_{\mathrm{St}})=(2-2g-p)G+\sum_a(h_a-s_a-k_a).}    \tag{E}
\end{equation}
\end{theorem}

\begin{proof}
The identity $d_1d_0=0$ is $d_\mu d_{\mathrm{act}}=0$.  Injectivity of
$d_{\mathrm{act}}$ leaves exactly the central $\CC^z$ in $H^0$.  The degree-one cohomology is the
displayed quotient $T$.  Surjectivity of $d_\mu$ leaves exactly the final central $\CC^z$ in
$H^2$.  Alternating dimensions and Theorem~\ref{thm:tangent-dimension} give
\eqref{eq:deformation-euler}.
\end{proof}

By Theorem~\ref{thm:microlocal-index}, the Stokes-root number $s_a$ is computed by an integer
directed microstalk sum.  Thus equations (M), (D), and (E) form a closed chain from coefficient
differences to vanishing quotients to tangent dimension; no separately quoted dimension polynomial
is used in the deduction.

\section{Type-$A$ marking legs}\label{sec:legs}

Attach a path of length $l_i\geq0$ to orbit vertex $i$.  Put
$y_{i,0}=x_i$, $y_{i,l_i+1}=0$, and
$d_{i,k}=y_{i,k}-y_{i,k+1}$ for $0\leq k\leq l_i$.  Then
$\sum_kd_{i,k}=x_i$.

\begin{lemma}[Path completion]\label{lem:path-completion}
If $d_0+\cdots+d_l=x$, then
\begin{equation}\label{eq:path-completion}
 \sum_{k=0}^{l}d_k^2=
 \sum_{k=0}^{l}\left(d_k-\frac{x}{l+1}\right)^2+\frac{x^2}{l+1}.
\end{equation}
\end{lemma}

\begin{proof}
Expand and use $\sum_kd_k=x$.
\end{proof}

\begin{theorem}[Leg-stable Witt law]\label{thm:legs}
Replacing every terminal square $x_i^2$ by its full type-$A$ path energy gives
\begin{align}\label{eq:leg-law}
 \widehat Q_A
 &=\sum_i\sum_{k=0}^{l_i}\left(d_{i,k}-\frac{x_i}{l_i+1}\right)^2
   +\sum_i\frac{x_i^2}{l_i+1}
   +\sum_e\ell_e\nu_e(x)^2-(H-1)M(x)^2.
\end{align}
Hence arbitrary finite marking legs do not create a second negative or null direction.
\end{theorem}

\begin{proof}
Apply Lemma~\ref{lem:path-completion} independently to every leg.  The first term is a positive
fluctuation complement.  The remaining positive matrix has diagonal
$\operatorname{diag}((l_i+1)^{-1})>0$ plus the cut Gram matrix.  It is therefore positive definite
and is perturbed by the same one rank square.  The proof of Theorem~\ref{thm:global-witt} applies.
\end{proof}

\section{Spectral consequence and sharpness}\label{sec:spectral}

Suppose an honest finite Stokes graph has symmetric adjacency matrix $B$ and its Cartan
polarization $C=2I-B$ is one of the forms covered by Theorem~\ref{thm:global-witt}, including its
type-$A$ legs.  Let
$\lambda_1(B)\geq\cdots\geq\lambda_n(B)$.

\begin{theorem}[Second-eigenvalue obstruction]\label{thm:lambda-two}
If $n\geq2$, then
\begin{equation}\label{eq:lambda-two}
                              \lambda_2(B)\leq2.
\end{equation}
If $\lambda_1(B)=2$, that eigenvalue is simple.  If $\lambda_1(B)>2$, exactly one adjacency
eigenvalue exceeds $2$ and none equals $2$.
\end{theorem}

\begin{proof}
The eigenvalues of $C$ are $2-\lambda_i(B)$.  Apply
\eqref{eq:global-inertia}.
\end{proof}

\begin{example}[Sharp Lorentz phase]\label{ex:triangle}
The symmetric matrix
\[
 B=\begin{pmatrix}0&6&4\\6&0&3\\4&3&0\end{pmatrix},\qquad
 C=2I-B,
\]
has $\one^TC\one=-20$, $\det C=-258$, and $\operatorname{tr}C=6$.  Its inertia is therefore
$(2,1,0)$.  The universal statement cannot be strengthened to positivity; one negative direction
is sharp.
\end{example}

\section{Novelty boundary and falsifiers}\label{sec:novelty}

Several ingredients have strong antecedents.  Finite Stokes filtrations and their categorical
forms are classical \cite{SabbahStokes}; irregular Riemann--Hilbert theory supplies analytic
realizations \cite{DKRH}; ultrametric matrices admit tree descriptions and inverse
$M$-matrix theory \cite{Martinez1994,McDonald1995,Nabben1995}; fission and wild character spaces
are developed in \cite{Boalch2014}; and Hiroe's theorem gives a scalar spectral-curve Euler
identity \cite{Hiroe2019}.

The new claim is their mixed dependency-closed polarization:
\begin{equation*}
 \text{Puiseux coefficients}
 \Longrightarrow \text{canonical threshold cuts}
 \Longrightarrow \text{quotient microstalk masses}
 \Longrightarrow \text{mixed Hom index}
 \Longrightarrow \text{positive coarea plus one rank line}.
\end{equation*}
Its global content is not present in a scalar rigidity formula: every puncture and every handle
must share the same single rank port, forcing the Witt bound in arbitrary genus.  The global
multi-Rees theorem and exact corona--relative-vanishing sequence identify how all collateral
constructions preserve that port.

Any of the following would falsify the theorem as stated.
\begin{enumerate}[label=(\arabic*),leftmargin=2.2em]
\item an equivariant coefficient family violating Lemma~\ref{lem:three-point};
\item failure of \eqref{eq:coarea} for its finite threshold partitions;
\item a directed exponential germ whose quotient microstalk mass differs from its degree;
\item a nonzero compatible rank-zero charge of nonpositive fused energy;
\item a nonnegative Rees specialization with two independent negative directions;
\item a positive corona block for which \eqref{eq:schur-completion} or
\eqref{eq:corona-exact} fails.
\end{enumerate}
Each item is a direct finite test.  Conversely, the theorem is not an analytic existence criterion:
a quadratic form satisfying the Witt law need not arise from a meromorphic connection.

\section{Conclusion}

Universal law is the factorization of the full mixed Stokes index through a positive microlocal
coarea space and one physical-rank line.  Principal conormals, cellular characteristic-cycle
coefficients, vanishing quotients, Rees graded channels, corona residuals, marking legs, and local
terms in the deformation complex are positive coordinates of that same space.  Fusion and topology
cannot create another defect because they see local data only through common physical rank.

This explains, in one statement, local Lorentz bounds, affine simplicity, capacity walls, global
handle monotonicity, Rees stability, Schur order independence, and the adjacency obstruction
$\lambda_2\leq2$.  More importantly, it replaces an interface between coefficient contacts and
microlocal multiplicities by an exact categorical quotient calculation.

\begin{thebibliography}{99}

\bibitem{BezrukavnikovKapranov}
R.~Bezrukavnikov and M.~Kapranov,
\emph{Microlocal sheaves and quiver varieties},
Ann. Fac. Sci. Toulouse Math. \textbf{25} (2016), 473--516.

\bibitem{Boalch2014}
P.~Boalch,
\emph{Geometry and braiding of Stokes data; fission and wild character varieties},
Ann. of Math. (2) \textbf{179} (2014), 301--365.

\bibitem{DKRH}
A.~D'Agnolo and M.~Kashiwara,
\emph{Riemann--Hilbert correspondence for holonomic $\mathcal D$-modules},
Publ. Math. Inst. Hautes Etudes Sci. \textbf{123} (2016), 69--197.

\bibitem{DKVanishing}
A.~D'Agnolo and M.~Kashiwara,
\emph{Enhanced nearby and vanishing cycles in dimension one and Fourier transform},
Publ. Res. Inst. Math. Sci. \textbf{59} (2023), 543--570.

\bibitem{Hiroe2017}
K.~Hiroe,
\emph{Ramified irregular singularities of meromorphic connections and plane curve singularities},
Josai Math. Monogr. \textbf{10} (2017), 161--192.

\bibitem{Hiroe2019}
K.~Hiroe,
\emph{Index of rigidity of differential equations and Euler characteristic of their spectral
curves}, arXiv:1907.12846.

\bibitem{Martinez1994}
S.~Mart\'{i}nez, G.~Michon and J.~San Mart\'{i}n,
\emph{Inverse of strictly ultrametric matrices are of Stieltjes type},
SIAM J. Matrix Anal. Appl. \textbf{15} (1994), 98--106.

\bibitem{McDonald1995}
J.~J.~McDonald, M.~Neumann, H.~Schneider and M.~J.~Tsatsomeros,
\emph{Inverse $M$-matrix inequalities and generalized ultrametric matrices},
Linear Algebra Appl. \textbf{220} (1995), 321--341.

\bibitem{Nabben1995}
R.~Nabben and R.~S.~Varga,
\emph{Generalized ultrametric matrices---a class of inverse $M$-matrices},
Linear Algebra Appl. \textbf{220} (1995), 365--390.

\bibitem{SabbahStokes}
C.~Sabbah,
\emph{Introduction to Stokes Structures},
Lecture Notes in Mathematics 2060, Springer, 2013.

\end{thebibliography}

\end{document}

\section{Cellular Stokes categories, microsupport, and vanishing cycles}\label{sec:cellular}

Let $\mathsf C$ be the walking cospan with objects $L,W,R$ and nonidentity arrows
$L\to W\leftarrow R$.  Put
\[
                         \Sh_{\mathrm{germ}}=\operatorname{Fun}(\mathsf C,\Vect_{\QQ}).
\]
It is an abelian category.  For $S=L$ or $R$, let $\rho_S(F):F(S)\to F(W)$ be the corresponding
incidence map and define
\begin{equation}\label{eq:microstalk-functor}
                         \mu_S(F)=F(W)/\im\rho_S(F).
\end{equation}

\begin{proposition}[Microlocal quotient functor]\label{prop:microfunctor}
Equation \eqref{eq:microstalk-functor} defines a functor
$\mu_S:\Sh_{\mathrm{germ}}\to\Vect_{\QQ}$.  For every germ $F$,
\begin{equation}\label{eq:nearby-vanishing-exact}
       F(S)\xrightarrow{\rho_S(F)}F(W)\longrightarrow\mu_S(F)\longrightarrow0
\end{equation}
is exact.
\end{proposition}

\begin{proof}
A morphism $\alpha:F\to G$ is a natural transformation.  Naturality of the square at $S\to W$
gives $\alpha_W(\im\rho_S(F))\subseteq\im\rho_S(G)$, so $\alpha_W$ induces a unique map on
quotients.  Identity and composition descend, proving functoriality.  Exactness and surjectivity
of the last map are the defining properties of a quotient vector space.
\end{proof}

For a finite-dimensional vector space $V$ and $\varepsilon\in\{0,1\}$, define the exponential
germ $E(V,\varepsilon)$ by
\[
 E(W)=V,\qquad (E(L),E(R))=
 \begin{cases}(V,0),&\varepsilon=1,\\(0,V),&\varepsilon=0,\end{cases}
\]
with the evident identity and zero incidence maps.

\begin{corollary}[One oriented vanishing cycle]\label{cor:one-oriented}
If $V\ne0$, exactly one of $\mu_L(E(V,\varepsilon))$ and
$\mu_R(E(V,\varepsilon))$ is zero; the other is canonically $V$.
\end{corollary}

\begin{proof}
The cokernel of $V\xrightarrow{1}V$ is zero and the cokernel of $0\to V$ is $V$.
\end{proof}

For $k>0$, label the $2k$ walls of the degree-$k$ cellular Stokes circle by
\[
                         \mathsf W_k=\{0,\ldots,k-1\}\times\{+,-\}.
\]
Place one arc between consecutive walls.  The entrance quiver has one arrow from every arc to each
endpoint wall.  No two nonidentity arrows compose, so its path category has only identities and
these incidence arrows.  Define the global cellular Stokes category
\begin{equation}\label{eq:global-cellular-category}
           \Sh_{\mathrm{St}}(S^1;k)=
           \operatorname{Fun}(\operatorname{Path}(\mathsf W_k),\Vect_{\QQ}).
\end{equation}
It is abelian.  The alternating exponential object has wall stalk $V$, alternating arc stalks
$V,0,V,0,\ldots$, and identity or zero incidence maps.

\begin{definition}[Cellular microsupport and characteristic cycle]\label{def:cellular-cc}
An oriented conormal $(w,S)$ belongs to the cellular microsupport of $F$ when
$\mu_S(F|_w)\ne0$.  Its cellular characteristic-cycle coefficient is
$\dim\mu_S(F|_w)$.  The directed cycle retains the $k$ conormals carrying the ordered decay
orientation.
\end{definition}

\begin{theorem}[Exponential microsupport calculation]\label{thm:exponential-microsupport}
Let $V$ have dimension $d$.  For the alternating degree-$k$ exponential object:
\begin{enumerate}[label=(\roman*),leftmargin=2.2em]
\item over every wall exactly one oriented conormal lies in cellular microsupport;
\item every nonzero characteristic-cycle coefficient equals $d$;
\item the full characteristic-cycle mass is $2kd$;
\item the directed decay mass is $kd$.
\end{enumerate}
\end{theorem}

\begin{proof}
Every wall germ is $E(V,0)$ or $E(V,1)$, so (i) and (ii) follow from
Corollary~\ref{cor:one-oriented}.  There are $2k$ walls.  Exactly one wall in each labelled
antipodal pair carries the ordered decay orientation, giving $k$ directed walls.  Summation proves
(iii) and (iv).
\end{proof}

The preceding theorem is a calculation inside an abelian category, not a numerical convention:
each coefficient is the dimension of the quotient functor \eqref{eq:microstalk-functor}, and
\eqref{eq:nearby-vanishing-exact} is its nearby-to-vanishing sequence.

\section{Microlocal Hom index and Galois descent}\label{sec:microlocal}

For branches $a,b$, let $E_{ab}$ be the cellular exponential Hom object of degree $k(a,b)$.
If orbit $i$ carries a coefficient space $V_i$ of dimension $d_i$, tensor $E_{ab}$ by
$\Hom(V_{[a]},V_{[b]})$.  Its fiber dimension is $d_{[a]}d_{[b]}$.

\begin{definition}[Normalized microlocal Hom mass]\label{def:microlocal-mass}
For orbit classes $i,j$, define
\begin{equation}\label{eq:microlocal-mass}
 \mathfrak m_{ij}(d)=\frac1N
 \sum_{a\in i}\sum_{b\in j}
 \sum_{\xi\in\operatorname{Decay}(a,b)}
       \dim\mu_\xi(E_{ab}\otimes\Hom(V_i,V_j)).
\end{equation}
\end{definition}

\begin{theorem}[Directed microlocal Hom index]\label{thm:microlocal-index}
For all $i,j$,
\begin{equation}\label{eq:microlocal-index}
                       \mathfrak m_{ij}(d)=s_{ij}d_id_j.      \tag{M}
\end{equation}
In particular, with one-dimensional coefficient spaces, the normalized directed microstalk mass
is exactly mixed Hom irregularity.
\end{theorem}

\begin{proof}
Theorem~\ref{thm:exponential-microsupport} gives directed mass
$k(a,b)d_id_j$ for each ordered pair.  Sum and divide by $N$:
\[
 \mathfrak m_{ij}(d)=\sum_{a\in i,b\in j}\frac{k(a,b)}N d_id_j
 =s_{ij}d_id_j.
\]
Lemma~\ref{lem:integrality} shows that normalization introduces no fractional multiplicity.
\end{proof}

Define the one-dimensional categorical microlocal pairing
\begin{equation}\label{eq:categorical-pairing}
                       \langle x,x\rangle_\mu
                       =\sum_{i,j}\mathfrak m_{ij}(1)x_ix_j.
\end{equation}

\begin{theorem}[Categorical microlocal coarea law]\label{thm:categorical-coarea}
For every orbit charge $x$,
\begin{align}\label{eq:categorical-coarea}
 Q_A(x)
 &=\sum_i x_i^2+M(x)^2-\langle x,x\rangle_\mu\notag\\
 &=\sum_i x_i^2+\sum_e\ell_e\nu_e(x)^2-(H-1)M(x)^2.
\end{align}
Thus microsupport coefficients, directed vanishing-cycle ranks, mixed Hom irregularities, and
threshold coarea coefficients are four presentations of the same polarization.
\end{theorem}

\begin{proof}
The first equality is Definition~\ref{def:local-index} and
Theorem~\ref{thm:microlocal-index}; the second is Theorem~\ref{thm:local-witt}.
\end{proof}

\section{Exact transport and irregular Riemann--Hilbert scope}\label{sec:transport}

We isolate the categorical statement needed to pass between realizations.  Let $\mathcal A$ be an
abelian category and let
\[
                         \Phi:\mathcal A\overset\sim\longrightarrow\Sh_{\mathrm{germ}}
\]
be an exact equivalence.  Define transported nearby and vanishing functors by
\[
             \Psi_W=\operatorname{ev}_W\Phi,qquad
             \phi_S=\mu_S\Phi.
\]

\begin{theorem}[Exact realization transport]\label{thm:transport}
For every $A\in\mathcal A$, the sequence
\begin{equation}\label{eq:transport-exact}
 \operatorname{ev}_S\Phi(A)\longrightarrow\Psi_W(A)
 \longrightarrow\phi_S(A)\longrightarrow0
\end{equation}
is exact and functorial.  If $\Phi(A)$ is an exponential germ, precisely one oriented transported
vanishing object has full coefficient rank.  All microstalk dimensions and therefore every term in
\eqref{eq:categorical-coarea} are invariant under replacing $\Phi$ by an isomorphic exact
realization equivalence.
\end{theorem}

\begin{proof}
Apply Proposition~\ref{prop:microfunctor} to $\Phi(A)$.  Exactness and functoriality are preserved
under composition.  The exponential statement is Corollary~\ref{cor:one-oriented}.  A natural
isomorphism of realization functors induces isomorphisms on stalk maps and hence on their
cokernels, preserving dimensions.
\end{proof}

\begin{remark}[Analytic boundary]
Irregular Riemann--Hilbert equivalences identify analytic Stokes objects with suitable sheaf
categories \cite{SabbahStokes,DKRH}; enhanced nearby and vanishing cycles in dimension one are
developed in \cite{DKVanishing}.  Theorem~\ref{thm:transport} proves the complete implication used
here once an exact realization identifies the local incidence diagram.  We make no converse claim
that the quadratic form alone reconstructs an analytic meromorphic connection or an equivalence.
\end{remark}

\section{Principal conormals and the multi-Rees law}\label{sec:rees}

For every dendrogram block $A_e$, let
\[
                      \mathcal H_e=\{x:\nu_e(x)=0\}.
\]
Its constant principal conormal is $d\nu_e=u_e=(u_{e,i})_i$.  Put
\[
 P_A(x)=\sum_i x_i^2+\sum_e\ell_e\nu_e(x)^2.
\]

\begin{proposition}[Principal-conormal Hessian]\label{prop:hessian}
\begin{equation}\label{eq:hessian}
                   \frac12\Hess(P_A)=I+\sum_e\ell_e,d\nu_e\otimes d\nu_e.
\end{equation}
The sole subtractive Hessian term in $Q_A$ is
$-(H-1)dM\otimes dM$.
\end{proposition}

\begin{proof}
Differentiate the squares in Theorem~\ref{thm:local-witt}.
\end{proof}

For parameters $t_e\in\RR$, define
\begin{equation}\label{eq:local-rees}
 Q_{A,\boldsymbol t}(x)=\sum_i x_i^2+
       \sum_et_e\ell_e\nu_e(x)^2-(H-1)M(x)^2.
\end{equation}
The associated-graded coordinate of $e$ is
$z_e=\sqrt{\ell_e}\nu_e$, and $\sum_ez_e^2$ is the block coarea energy.  Unit specialization is
$Q_A$; setting $t_e=0$ deletes exactly channel $e$.

At every marked point choose independent parameters $t_{a,e}$ and define the global Rees form by
the same fusion rule \eqref{eq:fused-form}.

\begin{theorem}[Global multi-Rees Witt law]\label{thm:global-rees}
If every $t_{a,e}\geq0$, then the global Rees form is a positive-definite terminal-plus-graded
form plus the same single common-rank square as in Theorem~\ref{thm:global-witt}.  Consequently
\[
              \indm(Q_{X,D,\boldsymbol t})\leq1,
              \qquad Q_{X,D,\boldsymbol t}>0
              \quad\text{on }\ker M\setminus\{0\}.
\]
The common-rank coefficient is independent of all Rees parameters.
\end{theorem}

\begin{proof}
Replace each local positive form by
$\sum_i x_i^2+\sum_et_{a,e}\ell_{a,e}\nu_{a,e}^2$.  Nonnegative parameters preserve
nonnegativity of every graded term, while the terminal coordinate sum remains positive definite.
Fusion changes only the already common rank square.  The hyperplane argument in
Theorem~\ref{thm:global-witt} applies unchanged.
\end{proof}

For real $t>0$, define
\begin{equation}\label{eq:threshold-relation}
              a\sim_t b\quad\Longleftrightarrow\quad a=b\ \hbox{ or }\ h(a,b)<t.
\end{equation}

\begin{theorem}[Canonical coefficient dendrogram]\label{thm:dendrogram}
Every relation $\sim_t$ is an equivalence relation.  Its partitions $\Pcal_t$ coarsen as $t$
increases.  The finitely many distinct blocks, ordered by inclusion, form a canonical rooted
dendrogram.  For distinct $a,b$, the smallest block containing both is born at $h(a,b)$.
\end{theorem}

\begin{proof}
Only transitivity needs proof.  If $a\sim_t b$ and $b\sim_t c$, then
Lemma~\ref{lem:three-point} gives $h(a,c)<t$.  Monotonicity in $t$ is immediate.  Only values
$k/N$ with $0\leq k\leq K$ occur, so there are finitely many partitions.  Join every block to the
least strictly larger block containing it and suppress repeated unary vertices.  The defining
threshold shows that the first common block of $a,b$ occurs at $h(a,b)$.
\end{proof}

Choose $H>\max_{a,b}h(a,b)$ and list all threshold values together with $0,H$.  On every open
interval between consecutive values, $\Pcal_t$ is constant.  Define the finite partition energy
\begin{equation}\label{eq:partition-energy}
 E_H(y)=\int_0^H\sum_{C\in\Pcal_t}\left(\sum_{a\in C}y_a\right)^2dt,
\end{equation}
where the integral denotes the corresponding finite sum of interval length times the constant
integrand.

\begin{theorem}[Coefficient coarea identity]\label{thm:coarea}
For every $y:A\to\RR$,
\begin{equation}\label{eq:coarea}
 E_H(y)=H\left(\sum_ay_a\right)^2-\sum_{a,b}h(a,b)y_ay_b.    \tag{C}
\end{equation}
Equivalently, if $e$ runs through the non-root dendrogram blocks, $\ell_e>0$ is their persistence
length, and $y(A_e)=\sum_{a\in A_e}y_a$, then
\begin{equation}\label{eq:cut-energy}
                         E_H(y)=\sum_e\ell_e y(A_e)^2.
\end{equation}
Moreover $E_H(y)>0$ for $y\ne0$.
\end{theorem}

\begin{proof}
Expand \eqref{eq:partition-energy}.  The coefficient of $y_ay_b$ is the measure of thresholds
$t\in(0,H)$ for which $a\sim_t b$.  This interval has length $H-h(a,b)$, including $a=b$ since
$h(a,a)=0$.  Summing gives \eqref{eq:coarea}.  Grouping the step integral by persistent blocks
gives \eqref{eq:cut-energy}.  For sufficiently small positive $t$, all blocks are singletons;
hence a nonzero coordinate contributes on an interval of positive length.
\end{proof}

\section{Orbit descent and the local Witt law}\label{sec:orbit}

Let $I=A/\langle\sigma\rangle$ and write $[a]$ for the orbit of $a$.  Put
\[
 r_i=|i|,\qquad
 s_{ij}=\sum_{a\in i}\sum_{b\in j}h(a,b),\qquad
 M(x)=\sum_ir_ix_i.
\]

\begin{lemma}[Galois integrality]\label{lem:integrality}
Each $s_{ij}$ is a nonnegative integer.
\end{lemma}

\begin{proof}
Decompose $i\times j$ into cyclic orbits.  Let one such orbit have length $L$ and let the leading
degree of $q_b-q_a$ be $k$.  Its stabilizer has order $N/L$ and contains $\sigma^L$.  Equivariance
forces $\zeta^{-Lk}=1$, hence $N$ divides $Lk$.  This pair orbit contributes
$Lk/N\in\NN$ to $s_{ij}$.  Sum over pair orbits.
\end{proof}

\begin{definition}[Polarized local Stokes index]\label{def:local-index}
For $x:I\to\RR$, set
\begin{equation}\label{eq:local-index}
             Q_A(x)=\sum_i x_i^2+M(x)^2-\sum_{i,j}s_{ij}x_ix_j.
\end{equation}
\end{definition}

For a dendrogram block $A_e\subseteq A$, define its descended incidence vector and charge by
\begin{equation}\label{eq:descended-cut}
        u_{e,i}=|A_e\cap i|,\qquad
        \nu_e(x)=\sum_i u_{e,i}x_i=\sum_{a\in A_e}x_{[a]}.
\end{equation}

\begin{theorem}[Local Stokes--microlocal Witt law]\label{thm:local-witt}
For every orbit charge $x$,
\begin{equation}\label{eq:local-witt}
 \boxed{
 Q_A(x)=\sum_i x_i^2+\sum_e\ell_e\nu_e(x)^2-(H-1)M(x)^2.}   \tag{W}
\end{equation}
Consequently $Q_A$ is positive definite on $\ker M$, and $\indm(Q_A)\leq1$.
\end{theorem}

\begin{proof}
Lift $x$ to $y_a=x_{[a]}$.  Then
\[
 \sum_{i,j}s_{ij}x_ix_j=\sum_{a,b}h(a,b)y_ay_b,
 \qquad \sum_ay_a=M(x).
\]
Insert Theorem~\ref{thm:coarea} into Definition~\ref{def:local-index}.  This gives
\eqref{eq:local-witt}.  On $\ker M$, the last square vanishes and $\sum_i x_i^2$ is strictly
positive away from zero.  A two-dimensional negative subspace would meet the codimension-one
space $\ker M$ nontrivially, contradiction.
\end{proof}

\begin{remark}[No orbit-saturation hypothesis]
The blocks $A_e$ may split a Galois orbit.  Formula \eqref{eq:descended-cut}, rather than
$u_{e,i}=r_i\one_{A_e}(i)$, is what makes Theorem~\ref{thm:local-witt} valid in the ramified case.
When every block is orbit-saturated the older rooted-core cut formula is recovered.
\end{remark}

\section{The global fused Witt law}\label{sec:global}

Let $D=\{p_1,\ldots,p_m\}$ with $m\geq1$.  At $p_a$ choose a spectrum $A_a$, orbit space $I_a$,
local rank $M_a$, and local index $Q_a$.  Define
\begin{equation}\label{eq:fibre-product}
 V_{X,D}=\{(x_a)\in\bigoplus_a\RR^{I_a}:M_1(x_1)=\cdots=M_m(x_m)\}.
\end{equation}
Write $M$ for common rank.  For genus $g$, define
\begin{equation}\label{eq:fused-form}
             Q_{X,D}(x)=\sum_{a=1}^{m}Q_a(x_a)+2(1-g-m)M^2.
\end{equation}

For each local spectrum choose $H_a$ as above and write
\[
 P_{X,D}(x)=\sum_a\left(\sum_{i\in I_a}x_{a,i}^2+
                    \sum_e\ell_{a,e}\nu_{a,e}(x_a)^2\right),
\]
\begin{equation}\label{eq:global-coefficient}
 \Gamma_{X,D}=2(1-g-m)-\sum_a(H_a-1).
\end{equation}

\begin{theorem}[Global coefficient-generated Witt law]\label{thm:global-witt}
On $V_{X,D}$,
\begin{equation}\label{eq:global-decomposition}
                  \boxed{Q_{X,D}=P_{X,D}+\Gamma_{X,D}M^2.}   \tag{GW}
\end{equation}
The form $P_{X,D}$ is positive definite.  Therefore
\begin{equation}\label{eq:global-inertia}
 Q_{X,D}>0\quad\text{on }\ker M\setminus\{0\},\qquad
 \indm(Q_{X,D})+\dim\ker(Q_{X,D})\leq1.                    \tag{I}
\end{equation}
\end{theorem}

\begin{proof}
Sum Theorem~\ref{thm:local-witt} and collect the common-rank squares, including the fusion term.
This proves \eqref{eq:global-decomposition}.  Positivity of $P_{X,D}$ follows from its terminal
sum $\sum_{a,i}x_{a,i}^2$.

For the inertia statement, $Q_{X,D}$ is positive on $\ker M\setminus\{0\}$.  Any two-dimensional
negative subspace meets $\ker M$ nontrivially, so the negative index is at most one.  The radical
meets $\ker M$ trivially, so restriction of $M$ to the radical is injective and the nullity is at
most one.  If a negative vector and a nonzero radical vector both existed, their span would contain
a nonzero rank-zero vector; because the radical is orthogonal to every vector, that vector would
have negative energy, contradiction.  Thus the sum in \eqref{eq:global-inertia} is at most one.
\end{proof}

\begin{corollary}[Handle and sewing laws]\label{cor:handle}
Increasing genus by one subtracts $2M^2$ and changes no local channel.  If two marked curves are
joined away from marked points, then their common-rank coefficients add and then decrease by $2$.
\end{corollary}

\begin{proof}
Both statements are immediate from \eqref{eq:global-coefficient} and
$\chi(X\#Y)=\chi(X)+\chi(Y)-2$.
\end{proof}
